% GENERATED FILE. Edit canonical lessons, then run npm run generate:books.
% canonical-source-hash: fnv1a64:389c87881a69bc1a
% canonical-lessons: FR-C347-puisque, FR-C347-etant-donne-que, FR-C347-vu-que, FR-C347-par-consequent, FR-C347-c-est-pourquoi

\chapter{Puisque et vu que}
\label{ch:fr-puisque-et-vu-que}
\hlchaptermodality{347}

\begin{chapteropening}
\textbf{By the end of this chapter:} \emph{I can give a known reason with since, given that and seeing that, and draw a result with consequently and that is why.}

Complete the last lesson of chapter 347: I can give a known reason with since, given that and seeing that, and draw a result with consequently and that is why.
\end{chapteropening}

\section[puisque]{puisque — since, as (a known reason)}

\label{lesson:FR-C347-puisque}



\begin{quote}
Before the new one: say the French for to lead to, then the French for to follow.
\end{quote}

\subsection*{You'll want to know: puisque}
\textbf{puisque} — \textquotedblleft{}since, as (a known reason)\textquotedblright{}.

\begin{quote}
Puisque tu es là, tu peux m'aider.
\end{quote}

\begin{quote}
\emph{Since you are here, you can help me.}
\end{quote}

Takes the indicative: \emph{puisque tu es là}. The reason is one both people already know. Before \emph{il} or \emph{elle}: \emph{puisqu'il}.

\begin{cousinweb}
\emph{puisque} is built from \emph{puis} \textquotedblleft{}then, afterwards\textquotedblright{} and \emph{que}. \emph{puis} is usually traced to a Vulgar Latin form of \emph{post} \textquotedblleft{}after\textquotedblright{} (Lewis \& Short).
\end{cousinweb}

\begin{grammarlens}[title={What you've built}]
One more word for telling, explaining and describing at B1.
\end{grammarlens}

\subsection*{Your turn}
\begin{itemize}
\raggedright
  \item \emph{Say it:} \emph{puisque}
  \item \emph{Say it:} \emph{puisque}, once more
  \item \emph{Say it:} a sentence of your own with \emph{puisque}
  \item \emph{Recall:} say the French for to lead to, then the French for to follow, then say \emph{puisque} again
\end{itemize}

\begin{tcolorbox}[breakable,colback=teal!4,colframe=teal!35!black,title={Before you move on}]
What does puisque mean? (\textquotedblleft{}Since, as (a known reason)\textquotedblright{}.) Say it once more.
\end{tcolorbox}
\section[étant donné que]{étant donné que — given that}

\label{lesson:FR-C347-etant-donne-que}



\begin{quote}
Before the new one: say the French for to follow, then the French for since.
\end{quote}

\subsection*{You'll want to know: étant donné que}
\textbf{étant donné que} — \textquotedblleft{}given that\textquotedblright{}.

\begin{quote}
Étant donné que le magasin est fermé, nous rentrons.
\end{quote}

\begin{quote}
\emph{Given that the shop is closed, we are going home.}
\end{quote}

Takes the indicative: \emph{étant donné que le magasin est fermé}. Careful style, frequent in writing.

\begin{cousinweb}
\emph{étant donné que} is built from \emph{étant}, the present participle of \emph{être}, and \emph{donné} \textquotedblleft{}given\textquotedblright{}. \emph{donner} is inherited from Latin \emph{dōnāre} \textquotedblleft{}to give as a present\textquotedblright{} (Lewis \& Short).
\end{cousinweb}

\begin{grammarlens}[title={What you've built}]
One more word for telling, explaining and describing at B1.
\end{grammarlens}

\subsection*{Your turn}
\begin{itemize}
\raggedright
  \item \emph{Say it:} \emph{étant donné que}
  \item \emph{Say it:} \emph{étant donné que}, once more
  \item \emph{Say it:} a sentence of your own with \emph{étant donné que}
  \item \emph{Recall:} say the French for to follow, then the French for since, then say \emph{étant donné que} again
\end{itemize}

\begin{tcolorbox}[breakable,colback=teal!4,colframe=teal!35!black,title={Before you move on}]
What does étant donné que mean? (\textquotedblleft{}Given that\textquotedblright{}.) Say it once more.
\end{tcolorbox}
\section[vu que]{vu que — seeing that}

\label{lesson:FR-C347-vu-que}



\begin{quote}
Before the new one: say the French for since, then the French for given that.
\end{quote}

\subsection*{You'll want to know: vu que}
\textbf{vu que} — \textquotedblleft{}seeing that\textquotedblright{}.

\begin{quote}
Vu qu'il pleut, nous restons à la maison.
\end{quote}

\begin{quote}
\emph{Seeing that it is raining, we are staying at home.}
\end{quote}

Takes the indicative: \emph{vu qu'il pleut}. A little more relaxed than \emph{étant donné que}.

\begin{cousinweb}
\emph{vu que} is built from \emph{vu}, the past participle of \emph{voir} \textquotedblleft{}to see\textquotedblright{}, which is inherited from Latin \emph{vidēre} \textquotedblleft{}to see\textquotedblright{} (Lewis \& Short).
\end{cousinweb}

\begin{grammarlens}[title={What you've built}]
One more word for telling, explaining and describing at B1.
\end{grammarlens}

\subsection*{Your turn}
\begin{itemize}
\raggedright
  \item \emph{Say it:} \emph{vu que}
  \item \emph{Say it:} \emph{vu que}, once more
  \item \emph{Say it:} a sentence of your own with \emph{vu que}
  \item \emph{Recall:} say the French for since, then the French for given that, then say \emph{vu que} again
\end{itemize}

\begin{tcolorbox}[breakable,colback=teal!4,colframe=teal!35!black,title={Before you move on}]
What does vu que mean? (\textquotedblleft{}Seeing that\textquotedblright{}.) Say it once more.
\end{tcolorbox}
\section[par conséquent]{par conséquent — consequently, therefore}

\label{lesson:FR-C347-par-consequent}



\begin{quote}
Before the new one: say the French for given that, then the French for seeing that.
\end{quote}

\subsection*{You'll want to know: par conséquent}
\textbf{par conséquent} — \textquotedblleft{}consequently, therefore\textquotedblright{}.

\begin{quote}
Il est malade ; par conséquent, il ne vient pas.
\end{quote}

\begin{quote}
\emph{He is ill; consequently, he is not coming.}
\end{quote}

Written style; in speech, \emph{donc} or \emph{alors}.

\begin{cousinweb}
\emph{conséquent} is a learned borrowing from Latin \emph{cōnsequēns} \textquotedblleft{}following, consistent\textquotedblright{} (EtymDB), the participle of \emph{cōnsequī} \textquotedblleft{}to follow after\textquotedblright{}, built on \emph{sequī} \textquotedblleft{}to follow\textquotedblright{} (Lewis \& Short).
\end{cousinweb}

\begin{grammarlens}[title={What you've built}]
One more word for telling, explaining and describing at B1.
\end{grammarlens}

\subsection*{Your turn}
\begin{itemize}
\raggedright
  \item \emph{Say it:} \emph{par conséquent}
  \item \emph{Say it:} \emph{par conséquent}, once more
  \item \emph{Say it:} a sentence of your own with \emph{par conséquent}
  \item \emph{Recall:} say the French for given that, then the French for seeing that, then say \emph{par conséquent} again
\end{itemize}

\begin{tcolorbox}[breakable,colback=teal!4,colframe=teal!35!black,title={Before you move on}]
What does par conséquent mean? (\textquotedblleft{}Consequently, therefore\textquotedblright{}.) Say it once more.
\end{tcolorbox}
\section[c'est pourquoi]{c'est pourquoi — that is why}

\label{lesson:FR-C347-c-est-pourquoi}



\begin{quote}
Before the new one: say the French for seeing that, then the French for consequently.
\end{quote}

\subsection*{You'll want to know: c'est pourquoi}
\textbf{c'est pourquoi} — \textquotedblleft{}that is why\textquotedblright{}.

\begin{quote}
Je n'ai pas d'argent ; c'est pourquoi je reste ici.
\end{quote}

\begin{quote}
\emph{I have no money; that is why I am staying here.}
\end{quote}

It brings in the result, after the cause.

\begin{cousinweb}
\emph{pourquoi} is built from \emph{pour} \textquotedblleft{}for\textquotedblright{}, from Latin \emph{prō}, and \emph{quoi} \textquotedblleft{}what\textquotedblright{}, from Latin \emph{quid} (Lewis \& Short): for what.
\end{cousinweb}

\begin{grammarlens}[title={What you've built}]
One more word for telling, explaining and describing at B1.
\end{grammarlens}

\subsection*{Your turn}
\begin{itemize}
\raggedright
  \item \emph{Say it:} \emph{c'est pourquoi}
  \item \emph{Say it:} \emph{c'est pourquoi}, once more
  \item \emph{Say it:} a sentence of your own with \emph{c'est pourquoi}
  \item \emph{Recall:} say the French for seeing that, then the French for consequently, then say \emph{c'est pourquoi} again
\end{itemize}

\begin{tcolorbox}[breakable,colback=teal!4,colframe=teal!35!black,title={Before you move on}]
What does c'est pourquoi mean? (\textquotedblleft{}That is why\textquotedblright{}.) Say it once more.
\end{tcolorbox}
